% Business-Paper auf Basis der Harwness-Vorlage.
\documentclass[11pt,a4paper]{scrartcl}
\newcommand{\harwLanguage}{german}
\newcommand{\harwAccentHex}{714B67}
\newcommand{\harwWarnHex}{C0392B}
\newcommand{\harwMainFont}{DejaVu Serif}
\newcommand{\harwSansFont}{DejaVu Sans}
\newcommand{\harwMonoFont}{DejaVu Sans Mono}
\usepackage{harw-report}
\title{Veröffentlichung von harw – Business-Paper zur Release-Entscheidung}
\subtitle{Release-Entscheidung}
\author{Mia}
\date{2026-10-02}
\begin{document}
\maketitle

\begin{merke}[Kernaussage]
Eine breite Veröffentlichung ist anhand der derzeit belegten Informationen nicht freizugeben. Bereite stattdessen einen begrenzten, beaufsichtigten Linux-Piloten vor und starte ihn erst, wenn Lizenz, Tests/CI, Artefakte, Installation und Rücknahme nachweislich geprüft sind; andernfalls verschieben.
\end{merke}
\begin{achtung}[Wichtig vorab]
Das Matrix-Game ist nach zwei von drei Runden ohne auswertbare Spielzüge beendet worden. Das ist ein Methoden-/Simulationsversagen, kein Release-Votum. Die folgende Business-Analyse ist davon getrennt und ihre Akteursperspektiven sind hypothetisch.
\end{achtung}
\tableofcontents
\newpage

\section{Kontext und Entscheidungsfrage}\label{sec:kontext}
Harwness beschreibt harw als Rust-Workspace für beaufsichtigte KI-Agenten mit expliziten Berechtigungen, typisierten Werkzeugen, dauerhaftem Ausführungszustand und optionaler Host-Sicherheitsdurchsetzung. Das README nennt den Entwicklungsstatus aktiv und weist darauf hin, dass Schnittstellen und Konfiguration zwischen Revisionen wechseln können (README.md:5–9). Die Workspace-Metadaten nennen Version 0.9.0 und `publish = false` (Cargo.toml:183–192).

\section{Empfehlung und Entscheidungslogik}\label{sec:empfehlung}
Die Veröffentlichung jetzt nicht als breit verfügbaren, produktionsreifen Release freigeben. Bevorzugt wird ein begrenzter und rücknehmbarer Pilot, jedoch ausschließlich nach Schließung der Mindest-Gates. Ein offenes kritisches Gate oder ein fehlender Rücknahmeweg bedeutet Verschiebung. Diese Empfehlung folgt aus dem Reife- und Nachweisstand, nicht aus einem Matrix-Game-Ergebnis.

\begin{figure}[htbp]
\centering
\begin{tikzpicture}[node distance=1.1cm]
\node[harwbox,fill=harwaccent!15,text width=5cm] (decision) {Entscheidung anhand überprüfbarer Mindest-Gates};
\node[harwbox,fill=harwgrau,below left=of decision] (delay) {Kritisches Gate offen: verschieben};
\node[harwbox,fill=harwgrau,below right=of decision] (pilot) {Mindest-Gates erfüllt: beaufsichtigter Pilot};
\draw[harwpfeil] (decision) -- (delay);
\draw[harwpfeil] (decision) -- (pilot);
\end{tikzpicture}
\caption{Analytische Entscheidungslogik, kein Simulationsergebnis.}\label{fig:logik}
\end{figure}

\section{Belegte Fakten und offene Annahmen}\label{sec:fakten}
\begin{tabularx}{\textwidth}{@{}lL@{}}
\toprule
\textbf{Kategorie} & \textbf{Aussage und Beleg} \\
\midrule
Projektfakt & README beschreibt frühen Stand 0.9.0, veränderliche Schnittstellen/Konfiguration und beaufsichtigte lokale Nutzung (README.md:5--7). \\
Projektfakt & README dokumentiert getaggte GitHub-Releases und nennt Linux-Archive für x86\_64/aarch64, Prüfsummenprüfung und Voraussetzungen bubblewrap/util-linux (README.md:63--80); diese Dokumentation belegt kein aktuell verfügbares konkretes Release-Artefakt. Quellinstallation ist beschrieben (Zeilen 82--91). \\
Projektfakt & Workspace-Version 0.3.0, \texttt{publish = false} (Cargo.toml:98--105). \\
Projektfakt & Changelog führt Änderungen unter Unreleased und Matrix-Game-Arbeit (CHANGELOG.md:7, 52--63). \\
Projektfakt & Packaging empfiehlt heute Source-Installation; Homebrew-Datei ist Formelvorlage mit TODO für URL, SHA-256 und Version (packaging/README.md:3--22). \\
Auftraggeberhinweis & Kimi OCR TOML-Aktivierbarkeit ist laut Mia in Arbeit, nicht fertig und nicht standardmäßig aktiv. \\
Projektfakt & Cargo.toml deklariert MIT OR Apache-2.0 (Zeile 108); LICENSE-MIT und LICENSE-APACHE liegen im Repository-Root. Rechteprüfung sämtlicher Release-Inhalte und Freigabe für eine konkrete Verteilung sind damit nicht belegt. \\
\bottomrule
\end{tabularx}

\section{Das Matrix-Game liefert kein Entscheidungsergebnis}\label{sec:game}
Der Bericht nennt Szenario `harw-release-entscheidung`, Lauf `20260924T212906-1`, Ende nach zwei von drei Runden (Matrix-Bericht: Zeilen 5–10). Vier institutionelle Perspektiven waren vorgesehen (Zeilen 20–60). In beiden Runden passten alle Sitze; das Rundenprotokoll führt acht Entscheidungen als `Forfeited` auf (Zeilen 68–104). Die Ergebnistabelle weist pro Akteur zwei verworfene/gepasste Züge aus; zusätzlich steht in „Offene Fragen“ die Angabe 16 gepasster Züge, die mit dieser Tabelle nicht übereinstimmt. Der Bericht nennt null Erfolge, null Misserfolge und eine offene Entscheidung (Zeilen 106–124, 162–164).

Es gibt daraus keine simulierte Zustimmung, Ablehnung oder belastbare Akteurspräferenz. Fehlender Debrief und fehlende Umpire-Synthese bestätigen, dass der Lauf keine Entscheidungsgrundlage geschaffen hat (Zeilen 127--159). Ein geänderter Neustart wurde nicht durchgeführt; es liegt kein auswertbares Matrix-Game-Ergebnis vor. Es werden weder ein weiterer Lauf noch fingierte Ergebnisse behauptet.

\section{Hypothetische Perspektiven und Szenariopfade}\label{sec:pfade}
Die folgenden Argumente sind analytische Hypothesen, keine Spielzüge oder tatsächlichen Positionen.

\begin{small}
\setlength{\tabcolsep}{3pt}
\noindent
\begin{tabularx}{\textwidth}{@{}>{\raggedright\arraybackslash}p{2.7cm}LL@{}}
\toprule
\textbf{Pfad} & \textbf{Hypothetische Argumente} & \textbf{Bedingung / Kehrseite} \\
\midrule
Pilot & Release-Team kann begrenzte Reichweite und transparente Grenzen vertreten. Erstanwendende können beaufsichtigte, rücknehmbare Erprobung bevorzugen. Integrationsperspektive kann einen kleinen Kreis mit klaren Installationsschritten unterstützen. Review-Perspektive kann einen Pilot als Informationsgewinn gegenüber einer breiten Zusage bewerten. & Nur nach nachweisbaren Mindest-Gates; klarer Pilotumfang, Supportgrenzen, Rücknahmeweg, keine Behauptung unbelegter Sicherheit oder Einsatzreife. \\
Breite Veröffent\-lichung & Reichweite und einfache Auffindbarkeit anstreben. Anwender könnten ein fertiges Einstiegserlebnis erwarten; Integrations- und Review-Perspektiven würden konsistente Pakete, Kompatibilität und belastbare Betriebs-/Supportbelege verlangen. & Derzeit nicht begründbar: Reifegradhinweis und Änderlichkeit sind dokumentiert, während Rechteprüfung sämtlicher Release-Inhalte, Freigabe der konkreten Verteilung, Test- und Artefaktnachweise hier offen sind. OCR nicht als fertige Funktion bewerben. \\
Verschieben & Release-Team kann die Entscheidung bis zur Schließung der Nachweislücken vertagen. Anwender und Integratoren können klarere Grenzen und reproduzierbare Installation verlangen. Review kann Nicht-Veröffentlichung als legitime Alternative werten. & Verzögerung schützt vor überzogenen Zusagen, verschiebt aber auch Feedback; kein Markt- oder Umsatzschaden wird mangels Daten beziffert. Kriterien und nächste Prüfung festlegen. \\
\bottomrule
\end{tabularx}
\end{small}

\section{Veröffentlichungsgates}\label{sec:gates}
\begin{enumerate}
\item \textbf{Recht:} Cargo.toml deklariert MIT OR Apache-2.0 und beide Lizenzdateien liegen im Repository-Root. Rechteprüfung sämtlicher Release-Inhalte und Freigabe der konkreten Verteilung sind damit nicht belegt; beides vorab prüfen und dokumentieren.
\item \textbf{Qualität:} Relevante Tests und CI für den vorgesehenen Release-Commit ausgeführt und dokumentiert; dieses Papier behauptet keinen Testlauf.
\item \textbf{Lieferumfang:} Konkrete Artefakte, Zielplattformen, Version und Prüfsummen erzeugt und verifiziert. README-Text allein ist kein Artefaktnachweis.
\item \textbf{Installation/Rücknahme:} Vorgesehener Installationspfad sowie Upgrade, Deinstallation und Rollback praktisch geprüft.
\item \textbf{Kommunikation:} Früher Entwicklungsstand und Grenzen offenlegen; Kimi OCR nicht als fertig oder standardmäßig aktiv darstellen.
\item \textbf{Betrieb:} Verantwortlichkeit, Supportumfang, Meldungsweg sowie Fortsetzungs-/Stoppkriterien benennen.
\end{enumerate}

\section{Risiken und Mitigationen}\label{sec:risiken}
\small
\begin{longtable}{@{}>{\raggedright}p{3.1cm}>{\raggedright}p{4.6cm}>{\raggedright\arraybackslash}p{7.3cm}@{}}
\toprule
\textbf{Risiko} & \textbf{Unsicherheit / Beleg} & \textbf{Mitigation} \\
\midrule\endhead
\bottomrule\endlastfoot
Rechte unklar & Rechteprüfung sämtlicher Release-Inhalte und Verteilung noch ungeprüft; MIT OR Apache-2.0 ist in Cargo.toml:108 deklariert. & Rechteprüfung und dokumentierte Lizenzentscheidung vor Verteilung. \\
Regressionen & Früher Stand und veränderliche Schnittstellen (README.md:7). & Release-Commit festschreiben, relevante Prüfungen dokumentieren, bekannte Grenzen offenlegen. \\
Paket-/Installations\allowbreak fehler & README beschreibt Linux-Pfade/Voraussetzungen; Packaging empfiehlt Source und Homebrew ist Vorlage. & Zielplattformen begrenzen; Installation, Prüfsummen und Rollback praktisch prüfen. \\
OCR-Erwartung & TOML-Aktivierung laut Mia in Arbeit. & Funktionsstatus präzise kommunizieren; nicht als Release-Zusage führen. \\
Supportlast & Kapazität nicht belegt. & Pilotumfang und Supportgrenzen festlegen; bei fehlender Kapazität verschieben. \\
Spiel\allowbreak fehlinterpretiert & 2/3 Runden, alle Sitze passen, kein auswertbarer Zug. & Simulationsversagen getrennt berichten, keine Spielentscheidung ableiten. \\
\end{longtable}
\normalsize

\section{Phasenplan und Entscheidungspunkt}\label{sec:plan}
\begin{tabularx}{\textwidth}{@{}>{\raggedright\arraybackslash}p{3.2cm}LL@{}}
\toprule
\textbf{Phase} & \textbf{Ergebnis} & \textbf{Entscheidungspunkt} \\
\midrule
Nachweis\-aufnahme & Lizenz, Release-Commit, CI/Tests, Artefakte, Plattformen, Installation und Rücknahme prüfen. & Offene kritische Punkte verhindern Start. \\
Release-Kandidat & Begrenzter Kandidat mit Dokumentation, Prüfsummen und Einschränkungen. & Review aller Mindest-Gates. \\
Beaufsichtigter Pilot & Nur nach Gate-Freigabe, mit begrenztem Umfang, Support und Rücknahme. & Auf Basis tatsächlicher Beobachtungen fortsetzen, nachbessern oder stoppen. \\
Breite Freigabe oder Verschiebung & Nachweise und Pilotbeobachtungen erneut prüfen. & Breite Veröffentlichung nur bei geschlossenen Gates; andernfalls vertagen. \\
\bottomrule
\end{tabularx}

\section{Bedingte Empfehlung und offene Punkte}\label{sec:schluss}
Die breite Veröffentlichung jetzt nicht freigeben. Einen begrenzten, beaufsichtigten Linux-Piloten vorbereiten und erst nach Erfüllung der Mindest-Gates starten. Sind Lizenz, Test-/CI-Nachweise, tatsächliche Artefakte oder Rücknahme ungeklärt, verschieben. Zeitangaben und Erfolgsschwellen bleiben offen, da keine belastbaren Aufwands- oder Nutzungsdaten vorliegen.

\textbf{Restrisiken:} Rechteprüfung, Teststatus, konkrete Artefakte und Rollback sind noch nachzuweisen; Supportkapazität und Kompatibilitätsumfang sind nicht quantifiziert. Kimi OCR ist laut Auftraggeberin in Arbeit. Die Spielauswertung ist nicht entscheidungsfähig. Autorin und Datum waren nicht angegeben und werden nicht ergänzt.

\appendix
\section{Belegverzeichnis}\label{sec:quellen}
\begin{itemize}
\item \path{README.md}: Zeilen 5--7, 13--29, 59--61, 63--95.
\item \path{Cargo.toml}: Zeile 108 (Lizenzdeklaration: MIT OR Apache-2.0).
\item \path{LICENSE-MIT} und \path{LICENSE-APACHE}: Lizenztexte im Repository-Root.
\item \path{CHANGELOG.md}: Zeile 7 und Zeilen 52--63.
\item \path{packaging/README.md}: Zeilen 3--13 und 15--22.
\item \path{matrix/harw-release-entscheidung-20260924T212906-1.md}: Zeilen 5--10, 14--18, 20--60, 68--104, 106--124, 127--170.
\end{itemize}
\end{document}
